\section{Discussion}
\label{sec:discussion}

\subsection{The Careful-but-Brief Trade-off}
\label{sec:tradeoff}

The most striking finding is that per-sentence deliberation improves the quality of individual sentences while reducing the number of sentences produced.
\iscotshort generates an average of 6.4 \texttt{<think>} blocks per response, each consuming tokens from a fixed budget of 1{,}024.
The result is a response with fewer but more carefully crafted sentences---higher lexical diversity, longer sentence length, and more frequent hedging---but less overall coverage of the topic.

This trade-off is fundamental, not incidental.
Any method that allocates part of a fixed token budget to reasoning necessarily reduces the space available for visible output.
The question is whether the per-sentence quality improvement is worth the cost in breadth.
For applications where every sentence must be reliable (\eg medical question answering, legal analysis), the answer is plausibly yes.
For applications where comprehensive coverage is more important (\eg tutorial generation, literature reviews), the trade-off may be unfavorable.

\subsection{Why LLM Judges Miss the Quality Shift}
\label{sec:judge_bias}

Our LLM-as-judge evaluation found that \standardcot wins 100\% of coherence and depth comparisons, despite producing text with the \emph{lowest} lexical diversity (TTR = 0.564) among all conditions.
This result is difficult to explain except as a length bias: \standardcot produces the longest responses (338.2 words on average), and the judge systematically equates length with quality.

This finding has implications beyond our study.
Any method that trades verbosity for precision---whether through inter-sentence deliberation, compression, or conciseness---will be penalized by length-biased evaluation.
Future work on sentence-level quality improvements should use length-controlled comparisons or human evaluation.

\subsection{When Does Per-Sentence Thinking Help?}
\label{sec:when_helps}

Our results suggest that \iscotshort is most beneficial for tasks involving potentially misleading questions.
On \truthfulqa, where questions are designed to elicit common misconceptions, \iscotshort achieves 98\% truthfulness vs.\ 95\% for \standard.
The 5 cases where \iscotshort succeeds and \standard fails involve questions with embedded false presuppositions---exactly the type of error that per-sentence deliberation is designed to catch.

On \gsmk, where the task is sequential numerical computation, \iscotshort provides no benefit.
Math reasoning requires a continuous chain of calculation rather than sentence-level reconsideration, so the forced pauses between sentences add overhead without improving the reasoning process.

\subsection{Limitations}
\label{sec:limitations}

\para{Single model.}
We test only \gptfour.
Other models may respond differently to inter-sentence CoT prompting, particularly models with different instruction-following capabilities or native thinking mechanisms (\eg Claude's extended thinking).

\para{Length confound.}
The shorter responses produced by \iscotshort make it difficult to disentangle the effects of deliberation from the effects of brevity.
A follow-up study should compare \iscotshort with an enlarged token budget (allowing comparable visible output length) against \standard.

\para{Sample size.}
With 100 \truthfulqa questions, the 3 percentage-point improvement in truthfulness (98\% vs.\ 95\%) does not reach statistical significance ($p = 0.453$).
A larger sample would be needed to confirm this effect.

\para{Automated text analysis.}
Our hedging and self-correction metrics rely on keyword lists, which may miss syntactic hedging or contextual self-correction that does not use our target phrases.

\para{Fixed token budget.}
The 1{,}024 max-token limit creates a hard trade-off between thinking and output that may not exist in settings with larger or unlimited budgets.
